\clearpage
\VerdierPhysicalPageAnchor{0027}
$\operatorname{Ext}_{\mathcal A}^{n}(X,Y)$ 的元素，即复形 $X$ 与 $Y$ 的
第 $n$ 个超 Ext（按 [1] 的意义）。若 $X$ 与 $Y$ 是有界复形（只有有限
多个分量非零），则从 $X$ 到 $Y$ 的次数为 $n$ 的同态，可通过推广 Yoneda
构造 [3] 的方法定义。$\mathbf D(\mathcal A)$ 中态射的复合正是 Yoneda
复合。

\VerdierSemanticAnchor{intro-6-plan-chapter-three-spectral-objects}
第 4 节研究通常的谱对象。$\mathcal A$ 的每个复形都产生三个取值于
$\mathbf D(\mathcal A)$ 的谱对象，从而自动写出基本的谱序列。此外，复形
及其态射的任意序列
$\cdots\to X_n\to X_{n+1}\to\cdots$（过渡态射不必是单态射）都会产生
一个谱对象。

\VerdierSemanticAnchor{intro-7}
\noindent\textbf{7.}\quad
总之，我们要为一种形式体系奠定基础，而这种工作包含它的一切不便：省略
三段论与连锁三段论；证明本身没有真正的数学困难，却常常需要一连串有时
（或者说总是）令人厌烦的核验。因此，只有当这一形式体系凭借其灵活性与
一般性，能够表述并证明“真正的”定理时，本工作才有价值。表述并证明上同调
对偶定理的需要，是这一形式体系最初的动力。已知在不同背景下存在若干形式
上极为相似的上同调对偶理论：

\VerdierSemanticAnchor{intro-7-duality-theories}
\begin{enumerate}
\renewcommand{\labelenumi}{\arabic{enumi})}
\VerdierSemanticAnchor{intro-7-duality-1}
\item 概形上凝聚层上同调的对偶理论 [4]。

\VerdierSemanticAnchor{intro-7-duality-2}
\item 概形 étale 拓扑中层上同调的对偶理论 [5], [6]。

这两个理论由 A. Grothendieck 创立。

\VerdierSemanticAnchor{intro-7-duality-3}
\item Galois 模上同调的对偶理论，由 J. Tate 与 J.-L. Verdier 创立 [7]。

\VerdierSemanticAnchor{intro-7-duality-4}
\item 局部紧空间上层上同调的对偶理论 [9], [10]。

\VerdierSemanticAnchor{intro-7-duality-5}
\item 凝聚解析层上同调的对偶理论。该理论目前仍基本上处于猜想阶段。

\VerdierSemanticAnchor{intro-7-duality-6}
\item 局部域、数域或函数域的对偶理论，以及它与类域论的关系。该理论尚待
发展；最初的步骤由 M. Artin 完成 [8]。
\end{enumerate}

\VerdierSemanticAnchor{intro-7-adjunction}
所有这些理论都以联系导出范畴之间某些函子的伴随性质为起点（或终点）。
